% 第8章 8.3.1：[0,1] 上 f_n(x)=x^n，n=1,2,4,8,16,32,64（颜色由浅到深）。函数族一致有界于 1，
% 逐点极限在 [0,1) 上为 0、在 x=1 处为 1（虚线与实心点），极限不连续，故没有一致收敛子列。解析式绘制。
\begin{tikzpicture}[font=\figurefont\footnotesize]
\pgfplotsset{every axis/.style={axis lines=left,
  x axis line style={draw=FigureStroke,line width=.6pt},y axis line style={draw=FigureStroke,line width=.6pt},
  tick style={draw=FigureStroke},tick label style={font=\figurefont\scriptsize,text=FigureInk},
  title style={font=\figurefont\footnotesize},label style={font=\figurefont\scriptsize},
  legend style={draw=none,fill=none,font=\figurefont\scriptsize,cells={anchor=west}}}}
\begin{axis}[width=64mm,height=42mm,xmin=0,xmax=1.05,ymin=0,ymax=1.1,xtick={0,0.5,1},ytick={0,0.5,1},xlabel={$x$},
  xlabel style={at={(axis description cs:1,0)},anchor=north west,yshift=1mm}]
\foreach \n/\c in {1/15,2/28,4/42,8/56,16/70,32/85,64/100}{
  \edef\tmp{\noexpand\addplot[draw=FigureBlue!\c!FigureBlueFill,line width=1pt,domain=0:1,samples=200] {x^\n};}\tmp}
\addplot[draw=FigureYellow,line width=1.1pt,dashed] coordinates {(0,0) (1,0)};
\addplot[only marks,mark=*,mark size=1.6pt,FigureYellow] coordinates {(1,1)};
\addplot[only marks,mark=*,mark size=1.6pt,mark options={fill=white,draw=FigureYellow}] coordinates {(1,0)};
\node[anchor=west,font=\figurefont\scriptsize] at (axis cs:0.12,0.22) {$n=1$};
\node[anchor=east,font=\figurefont\scriptsize] at (axis cs:0.9,0.75) {$n=64$};
\draw[FigureInk,line width=.4pt] (axis cs:0.9,0.75)--(axis cs:0.993,0.75);
\end{axis}
\end{tikzpicture}
